\documentclass[12pt]{article}
\usepackage[utf8]{inputenc}
\usepackage[table,xcdraw]{xcolor}
\usepackage{array}
\usepackage{makecell}
\usepackage{geometry}
\usepackage{graphicx}
\usepackage{multirow}
\usepackage{enumitem}
\usepackage{amsmath}
\usepackage{listings}
\usepackage{color}
\usepackage{hyperref}
\usepackage{spverbatim}
\usepackage[linesnumbered,ruled,vlined]{algorithm2e}
\usepackage[backend=bibtex]{biblatex}
\addbibresource{referencias.bib}
\usepackage{amssymb}
\usepackage{gensymb}
\usepackage{fancyhdr} % <- Nuevo paquete

\geometry{margin=1.7cm}

% IMPORTANTE: ajusta headheight para reservar espacio al encabezado
\setlength{\headheight}{3cm} % Ajusta según el alto real del encabezado
\setlength{\headsep}{0.5cm}  % Espacio entre encabezado y texto

\renewcommand{\arraystretch}{1.1}
\newcolumntype{P}[1]{>{\raggedright\arraybackslash}p{#1}}

% === Encabezado personalizado con imagen y tabla ===
\pagestyle{fancy}
\fancyhf{} % Limpiar encabezados/footers

\fancyhead[L]{%
  \begin{tabular}{|@{}p{2.5cm}@{}|p{9cm}|p{4.5cm}|}
    \hline
    \multirow{2}{*}{\includegraphics[height=0.6cm]{image.png}} & \textbf{VICERRECTORADO DOCENTE} & Código: GUIA-PRL-001 \\
    & \textbf{CONSEJO ACADÉMICO} & Aprobación: 2016/04/06 \\
    \hline
    \multicolumn{3}{|p{16cm}|}{Formato: Guía de Práctica de Laboratorio / Talleres / Centros de Simulación} \\
    \hline
  \end{tabular}
}

\fancyfoot[C]{} % Pie vacío

\begin{document}

% Elimina o reduce mucho el espacio vertical inicial
%\vspace{4.5cm} 

% ===== FORMATO DE PRÁCTICA (tabla principal) =====
\noindent
\begin{tabular}{|P{2cm} P{2.5cm} P{2.5cm} P{7cm} P{2cm} P{2cm}|}
\hline
\multicolumn{3}{|c|}{
  \parbox[c][2.2cm][c]{7cm}{\centering
    \includegraphics[height=1.8cm]{image.png}
  }
} 
& 
\multicolumn{3}{c|}{
  \parbox[c][2.2cm][c]{10cm}{\centering
  \textbf{FORMATO DE GUÍA DE PRÁCTICA DE LABORATORIO / TALLERES / CENTROS DE SIMULACIÓN –PARA DOCENTES}}
} \\ \hline

\multicolumn{3}{|P{7cm}|}{CARRERA: TELECOMUNICACIONES} & \multicolumn{3}{P{10cm}|}{ASIGNATURA: } \\ \hline
\multicolumn{2}{|P{4.5cm}|}{NRO. PRÁCTICA:} & \multicolumn{1}{P{2.5cm}|}{2} & \multicolumn{3}{P{10cm}|}{TÍTULO PRÁCTICA: Inicializacion en ESP32 manejo de ADC y GPIO } \\ \hline

% OBJETIVOS
\multicolumn{6}{|P{17cm}|}{
\begin{tabular}[c]{@{}P{16.8cm}@{}}
\textbf{OBJETIVO(S):}\\[2pt]
\begin{itemize}[leftmargin=*, nosep]
\item Programar el ESP32 WROOM-32 en MicroPython para realizar control digital básico mediante GPIO.

\item Implementar la lectura analógica de un potenciómetro utilizando el ADC interno del ESP32.

\item Comprobar la interacción entre hardware y software mediante pruebas prácticas con LED, pulsador y sensor analógico.
  

\end{itemize}
\end{tabular}
} \\ \hline









% INSTRUCCIONES
\multicolumn{2}{|P{4.5cm}|}{\multirow{3}{*}{INSTRUCCIONES}} & \multicolumn{4}{P{12.5cm}|}{1. Configura los pines del LED, pulsador, ADC e implementar un sensor DHT en un solo archivo \texttt{main.py}.} \\ \cline{3-6} 
\multicolumn{2}{|P{4.5cm}|}{} & \multicolumn{4}{P{12.5cm}|}{2. Implementa funciones separadas para parpadeo, lectura de botón, ADC y DHT11.} \\ \cline{3-6} 
\multicolumn{2}{|P{4.5cm}|}{} & \multicolumn{4}{P{12.5cm}|}{3. En el bucle principal, ejecuta cada tarea periódicamente e imprime su estado por UART con \texttt{print()}.} \\ \cline{3-6} 
\multicolumn{2}{|P{4.5cm}|}{} & \multicolumn{4}{P{12.5cm}|}{4. Guarda el archivo como \texttt{main.py} en el ESP32 y observa los mensajes en la consola de Thonny.} \\ \hline

\multicolumn{6}{|P{17cm}|}{\centering \textbf{ACTIVIDADES POR DESARROLLAR}} \\ \hline
\multicolumn{6}{|P{17cm}|}{1. Revisar el código presentado en la guía.} \\ \hline
\multicolumn{6}{|P{17cm}|}{2. Implementar los programas en un solo archivo, monitoreando la tarea ejectuada mediante comunicacion Serial.} \\ \hline
\multicolumn{6}{|P{17cm}|}{3. Presentar los resultados en el monitor serial.} \\ \hline


% RESULTADOS
\multicolumn{6}{|P{17cm}|}{
\begin{tabular}[c]{@{}P{16.8cm}@{}}
\textbf{RESULTADO(S) OBTENIDO(S):}\\[2pt]
\begin{itemize}[leftmargin=*, nosep]
  \item Realizar procesos de investigación la configuración de ESP32 mediante micropython.
  \item Entender las aplicaciones de los microcontroladores y microprocesadores para manejo de GPIO de propósito específico.
\end{itemize}
\end{tabular}
} \\ \hline

% CONCLUSIONES
\multicolumn{6}{|P{17cm}|}{
\begin{tabular}[c]{@{}P{16.8cm}@{}}
\textbf{CONCLUSIONES:}\\[2pt]
Aprenden a trabajar en grupo dentro de plazos de tiempo establecidos, manejando software de diseño, colocando las características y restricciones propuestas.
\end{tabular}
} \\ \hline

% RECOMENDACIONES
\multicolumn{6}{|P{17cm}|}{
\begin{tabular}[c]{@{}P{16.8cm}@{}}
\textbf{RECOMENDACIONES:}\\[2pt]
Realizar el trabajo dentro del tiempo establecido.
\end{tabular}
} \\ \hline

\end{tabular}

\vspace{1cm}

\noindent Docente / Técnico Docente: \rule{8cm}{0.4pt} \\[0.5cm]
\noindent Firma: \rule{9cm}{0.4pt}

\newpage



\begin{table}[h!]
\centering
\begin{tabular}{|P{2cm} P{2.5cm} P{2.5cm} P{10cm}|}
\hline
\multicolumn{3}{|c|}{
  \parbox[c][2.2cm][c]{7cm}{\centering
    \includegraphics[height=1.8cm]{image.png}
  }
} 
& 
\parbox[c][2.2cm][c]{10cm}{\centering
  \textbf{FORMATO DE GUÍA DE PRÁCTICA DE LABORATORIO / TALLERES / CENTROS DE SIMULACIÓN – PARA ESTUDIANTES}
} \\ \hline

\multicolumn{3}{|P{7cm}|}{\textbf{CARRERA:}} & \textbf{ASIGNATURA:} Sistemas Embebidos \\ \hline
\multicolumn{2}{|P{4.5cm}|}{\textbf{NRO. PRÁCTICA:2}} &  & \textbf{TÍTULO PRÁCTICA: Inicializacion en ESP32 manejo de ADC y GPIO} \\ \hline

\multicolumn{4}{|P{17cm}|}{
  \begin{minipage}[t]{16.8cm}
    \textbf{OBJETIVO(S) ALCANZADO(S):} \\[4pt]
    \begin{itemize}[leftmargin=*, nosep]
      \item 
      \item 
      \item 
    \end{itemize}
  \end{minipage}
} \\ \hline


\multicolumn{4}{|c|}{\textbf{ACTIVIDADES DESARROLLADAS}} \\ \hline
\multicolumn{4}{|P{17cm}|}{
  \begin{itemize}[leftmargin=*, nosep]
    \item[1.] 
    \item[2.] 
    \item[3.] 
    \item[4.] 
    \item[5.] 
    \item[6.] 
    \item[N.] 
  \end{itemize}
} \\ \hline

\multicolumn{4}{|P{17cm}|}{
  \parbox[t]{16.8cm}{
    \textbf{RESULTADO(S) OBTENIDO(S):} \\[2pt]
    \vspace{2cm}
  }
} \\ \hline


\multicolumn{4}{|P{17cm}|}{
  \parbox[t]{16.8cm}{
    \textbf{CONCLUSIONES:} \\[2pt]
    \vspace{2cm}
  }
} \\ \hline


\multicolumn{4}{|P{17cm}|}{
  \parbox[t]{16.8cm}{
    \textbf{RECOMENDACIONES:} \\[2pt]
    \vspace{2cm}
  }
} \\ \hline


\end{tabular}
\end{table}

\vspace{1cm}

\noindent Nombre del estudiante: \rule{8cm}{0.4pt} \\[0.5cm]




\end{document}
